\chapter{Valutazione dei controllori di concorrenza}
\label{cap:25-valutazione-concorrenza}

Questo capitolo riporta per esteso le campagne che hanno prodotto i risultati
citati nel \cref{cap:12-concurrency-control}. \`E il capitolo pi\`u ricco di
risultati negativi del documento, ed \`e organizzato in modo che ciascuno preceda
la correzione che ha motivato.

Tutte le misure sono di agosto 2026 e provengono dall'impianto descritto nel
\cref{cap:22-metodo-sperimentale}. Dove un'affermazione \`e stata fatta e poi
confutata dalla misura, il documento conserva la confutazione anzich\'e
l'affermazione.

\section{Il confronto fra i due controllori per funzione}

\subsection{Impianto}

Due funzioni su un solo control plane, vincolate agli stessi quattro core, coda
di 100 ciascuna, tetto per funzione 8, budget totale \code{BUDGETED} pari a 12.
Carico a ciclo chiuso con picco di 103 richieste offerte per funzione su tre
finestre: una funzione da sola, poi le due in antifase, poi le due in fase.
Entrambi i modi hanno ricevuto esattamente lo stesso carico.

\subsection{Risultati aggregati}

\begin{center}
\small
\begin{adjustbox}{max width=\textwidth}
\begin{tabular}{@{}lrr@{}}
\toprule
 & \textbf{\code{ADAPTIVE\_PER\_POD}} & \textbf{\code{BUDGETED}} \\
\midrule
richieste & 899\,684 & 893\,841 \\
rifiutate & 7 (0{,}001\%) & 33 (0{,}004\%) \\
p95 end-to-end & 80{,}2\,ms & 79{,}7\,ms \\
p99 end-to-end & 90{,}2\,ms & 87{,}2\,ms \\
\bottomrule
\end{tabular}
\end{adjustbox}
\end{center}

I due modi sono entro l'1\% l'uno dall'altro end-to-end. Preso da solo, questo
risultato direbbe che il modo a budget non porta alcun beneficio. \`E la
conclusione sbagliata, per due ragioni distinte che le sezioni seguenti
sviluppano.

\subsection{La concorrenza concessa, per finestra}

\begin{center}
\small
\begin{adjustbox}{max width=\textwidth}
\begin{tabular}{@{}lrrrrrr@{}}
\toprule
\textbf{Finestra} & \textbf{AD. java} & \textbf{AD. lite} & \textbf{somma} &
\textbf{BU. java} & \textbf{BU. lite} & \textbf{somma} \\
\midrule
una sola & 7{,}8 & 4{,}1 (idle) & --- & 8{,}0 & 1{,}2 (idle) & --- \\
antifase & 8{,}0 & 5{,}0 & 13{,}0 & 8{,}0 & 2{,}9 & \textbf{10{,}9} \\
in fase & 7{,}9 & 5{,}5 & 13{,}4 & 8{,}0 & 3{,}7 & \textbf{11{,}7} \\
\bottomrule
\end{tabular}
\end{adjustbox}
\end{center}

\code{BUDGETED} mantiene il totale sotto il proprio budget di 12 e sposta
capacit\`a fra le funzioni al muoversi del carico. \code{ADAPTIVE\_PER\_POD} non
ha alcun meccanismo per fare n\'e l'una n\'e l'altra cosa, perch\'e i suoi due
controllori indipendenti non sanno l'uno dell'esistenza dell'altro. \`E la
propriet\`a per cui il modo esiste, ed \`e visibile qui e non negli aggregati.

\subsection{L'asimmetria non \`e iniquit\`a}

Le due funzioni ricevono concorrenza molto diversa, e questo \emph{non} \`e un
difetto dell'allocatore. Sono runtime diversi: \code{lite} serve in 2{,}15\,ms
contro 5{,}57\,ms, quindi la legge di Little le assegna un limite pi\`u piccolo.
A 1\,493 richieste al secondo le servono circa 3{,}2 in volo, e oltre non compra
completamenti.

\code{BUDGETED} le ha concesso 2{,}6 in media contro i 4{,}9 di
\code{ADAPTIVE} \emph{a parit\`a di throughput}, e il suo tempo di servizio \`e
sceso del 37\%. Met\`a della concorrenza, lo stesso lavoro svolto.

\subsection{Un errore di calibrazione}

Il \cref{cap:22-metodo-sperimentale} riporta l'errore che ha invalidato la prima
esecuzione di questo confronto e la regola che ne \`e stata derivata. Va
richiamato qui perch\'e riguarda direttamente l'interpretabilit\`a della tabella
precedente: con budget 8 anzich\'e 12, la coda traboccava per aritmetica in un
modo e non nell'altro, producendo 86\,302 rifiuti che non dicevano nulla sui
controllori.

\section{La coda \`e dove sta la latenza}

\code{function\_latency\_ms} \`e misurata dal control plane fra dispatch e
completamento: \`e il round trip, non il tempo di esecuzione interno alla
funzione. Il governor decide da essa.

Sotto il profilo a raffica l'attesa media in coda \`e stata di \textbf{37--43\,ms}
contro un tempo di servizio di circa \textbf{5\,ms}. La grandezza da cui il
controllore si orienta \`e circa \textbf{un settimo} di ci\`o che il chiamante
sperimenta.

La conseguenza \`e misurabile, non teorica. In un confronto iniziale
\code{BUDGETED} ha ridotto il tempo di servizio del 31\% su una funzione e del
44\% sull'altra --- entrambe molto dentro un SLO di 10\,ms --- e ha
\textbf{peggiorato} il p95 end-to-end, perch\'e l'attesa creata dal suo stesso
limite ha pi\`u che assorbito il guadagno.

\begin{quote}
Un limite di concorrenza non rimuove lavoro: lo sposta nel buffer.
\end{quote}

\section{Il ciclo chiuso non poteva porre la domanda}

Prima di poter valutare un controllore che decide dal sojourn, \`e stato
necessario correggere l'impianto.

Sotto il profilo \code{burst} ogni utente virtuale tiene una richiesta, quindi
il numero nel sistema \`e fissato dal generatore e la profondit\`a di coda \`e
\code{VU} $-$ \code{limite}. Due conseguenze, entrambe misurate:

\begin{itemize}[leftmargin=*]
  \item Portare il limite da 4 a 8 ha accorciato di quattro una coda profonda 99.
  Nessun controllore pu\`o muovere l'attesa in modo apprezzabile.
  \item La coda si assesta a 95 su 100 per costruzione, permanentemente sopra
  qualunque soglia di guardia. La guardia di ammissione del modo \code{SOJOURN}
  scatta a ogni tick e il modello non parla mai: la riesecuzione legge
  \code{8 [8--8]} per tutta la sua durata.
\end{itemize}

\begin{quote}
\`E per questo che due controllori che leggevano segnali completamente diversi
sono finiti entro il 2\% l'uno dall'altro su questo profilo. La somiglianza era
una propriet\`a dell'impianto.
\end{quote}

\`E il risultato metodologico centrale del capitolo, e la ragione per cui il
confronto aggregato della prima sezione va letto con cautela: quel confronto
girava su un impianto incapace di distinguere.

\section{Decidere dalla latenza del chiamante: due fallimenti e un risultato}

\subsection{Fallimento 1: il gradiente sul sojourn}

Alimentare la regola a gradiente esistente con la latenza end-to-end \`e
instabile, e \`e stato scartato prima dell'implementazione. La regola riduce il
limite quando la latenza eccede l'obiettivo; un limite pi\`u piccolo drena la
coda pi\`u lentamente, quindi l'attesa cresce, quindi si riduce ancora, fino al
pavimento.

\subsection{Fallimento 2: la ricerca del minimo}

Il sojourn ha un minimo interno dove il tempo di servizio non ce l'ha, quindi la
ricerca \`e ben posta in linea di principio. \`E stata implementata e falsificata:

\begin{center}
\small
\begin{adjustbox}{max width=\textwidth}
\begin{tabular}{@{}lrr@{}}
\toprule
 & \textbf{\code{ADAPTIVE\_PER\_POD}} & \textbf{\code{SOJOURN}, a ricerca} \\
\midrule
rifiutate & 32 (0{,}004\%) & \textbf{36\,181 (4{,}275\%)} \\
limite [min--max] & 7{,}8 [4--8] & 4{,}0 [1--8] \\
\bottomrule
\end{tabular}
\end{adjustbox}
\end{center}

Il punto di lavoro scelto non era il problema: a limite 4 consegnava il 96{,}5\%
del throughput con met\`a del tempo di servizio. Il problema \`e che \emph{il
limite fa due lavori}: insieme alla coda \`e anche la finestra di ammissione,
perch\'e una funzione trattiene \code{limite + queueSize} e rifiuta il resto.
Dimensionarlo dalla sola latenza ha fatto perdere il 4\% del traffico.

La causa profonda \`e un problema di attribuzione: la ricerca non pu\`o
distinguere ``la mia mossa ha aiutato'' da ``il carico \`e calato''. In un
avvallamento ha attribuito il sollievo naturale alla propria decisione di
restringere, ed era ancora piccola quando il carico \`e tornato.

\subsection{Il risultato: calcolare da un modello}

Sotto arrivi aperti, con picco sopra la capacit\`a misurata, stesso carico,
stessi core, stesse code, entrambi i modi per funzione:

\begin{table}[htbp]
\centering
\caption{\code{ADAPTIVE\_PER\_POD} contro \code{SOJOURN} sotto arrivi aperti.}
\label{tab:sojourn}
\small
\begin{adjustbox}{max width=\textwidth}
\begin{tabular}{@{}lrrl@{}}
\toprule
 & \textbf{\code{ADAPTIVE\_PER\_POD}} & \textbf{\code{SOJOURN}} & \\
\midrule
offerte & 659\,973 & 659\,995 & identiche per costruzione \\
\textbf{servite} & 549\,856 & \textbf{567\,210} & \textbf{$+3{,}2\%$} \\
rifiutate & 110\,117 (16{,}7\%) & \textbf{92\,785 (14{,}1\%)} &
\textbf{$-15{,}7\%$} \\
\textbf{p95 end-to-end} & 120{,}48\,ms & \textbf{107{,}68\,ms} &
\textbf{$-10{,}6\%$} \\
\textbf{p99 end-to-end} & 157{,}67\,ms & \textbf{133{,}82\,ms} &
\textbf{$-15{,}1\%$} \\
attesa in coda & 35{,}61\,ms & 31{,}03\,ms & $-12{,}9\%$ \\
tempo di servizio & 1{,}69\,ms & 2{,}00\,ms & $+18\%$ \\
limite medio & 3{,}9 & 2{,}8 & \\
\bottomrule
\end{tabular}
\end{adjustbox}
\end{table}

Lo scambio \`e andato nella direzione prevista: il tempo di servizio peggiora del
18\%, l'attesa migliora del 13\%, e poich\'e l'attesa \`e venti volte il tempo di
servizio il chiamante ci guadagna. L'obiettivo --- stesso throughput, stessi
rifiuti, p95 end-to-end inferiore --- era stato scritto nello scenario
\emph{prima} dell'esecuzione, e tutti e tre sono stati superati.

Il dettaglio istruttivo \`e che \code{SOJOURN} ha ottenuto questo con
\emph{meno} concorrenza media. Non vince aggiungendo capacit\`a ma collocandola
nel tempo, che \`e ci\`o per cui la predizione era stata introdotta.

\subsection{Caveat}

Riportati dal progetto e da prendere sul serio:

\begin{itemize}[leftmargin=*]
  \item Una sola esecuzione per modo, non ripetuta. Differenze del 10--15\% sono
  grandi ma non stabilite.
  \item Entrambi i sistemi stanno scartando il 14--16\% del carico, perch\'e il
  picco a ciclo aperto eccede genuinamente la capacit\`a. Il confronto \`e fra
  due sistemi sovraccarichi, non fra due sistemi comodi.
  \item \code{dropped\_iterations} \`e stato 21 e 0, quindi il generatore non
  stava rinunciando e i rifiuti sono della piattaforma.
\end{itemize}

Il secondo caveat \`e il pi\`u limitante: \code{SOJOURN} non \`e mai stato
osservato scegliere un limite in condizioni comode.

\section{Co-tenancy: cross-talk senza contesa}

\subsection{Il fenomeno}

Due funzioni su un control plane, una a carico costante mentre l'altra va e
viene. Quando la vicina arriva, la funzione stabile perde il 28\% del proprio
throughput e guadagna il 36\% di latenza.

\subsection{L'ipotesi ovvia \`e sbagliata}

La spiegazione immediata era la contesa di CPU. \`E stato aggiunto
\code{nanofaas.container-local.cpuset} per vincolare entrambe le funzioni agli
stessi quattro core, e l'esecuzione \`e stata ripetuta:

\begin{center}
\small
\begin{adjustbox}{max width=\textwidth}
\begin{tabular}{@{}lrr@{}}
\toprule
 & \textbf{4 core propri ciascuna} & \textbf{core 0--3 condivisi} \\
\midrule
java da sola & 1\,762 rps / 4{,}57\,ms & 1\,762 rps / 4{,}57\,ms \\
java con la vicina & 1\,274 rps / 6{,}23\,ms & 1\,274 rps / 6{,}23\,ms \\
combinato durante la sovrapposizione & 2\,769 rps & 2\,751 rps \\
\bottomrule
\end{tabular}
\end{adjustbox}
\end{center}

Identiche a tre cifre significative. Vincolare entrambe le funzioni a quattro
core condivisi costa lo 0{,}6\% del throughput combinato, che \`e rumore.

I quattro core non erano mai il vincolo attivo: a circa 2\,760 richieste al
secondo ciascuna richiesta pu\`o costare al massimo 1{,}45\,ms di CPU, e
l'obiettivo \`e raggiunto. Il cross-talk viene da ci\`o che le due funzioni
condividono in \emph{entrambe} le configurazioni --- l'unico control plane ---
non dai core.

\subsection{Perch\'e il \code{cpuset} resta}

Pur non essendo la spiegazione cercata, la funzionalit\`a \`e stata mantenuta:
rende il budget di concorrenza una divisione di una capacit\`a che
\emph{esiste}. Senza di essa, su una macchina con core in eccesso, la premessa
del modo \code{BUDGETED} \`e falsa e il modo non pu\`o essere valutato
(\cref{cap:15-deployment-provider}).

\begin{damisurare}
Il meccanismo esatto del cross-talk attraverso il control plane condiviso non \`e
stato isolato. I candidati --- contesa sul thread dello scheduler, sull'event
loop, sulle strutture concorrenti delle code --- sono distinguibili con la
profilazione, e non sono stati distinti.
\todomis{origine del cross-talk fra funzioni co-residenti nel control plane}
\end{damisurare}

\section{Un confondente scoperto: i due \code{word-stats}}

Registrato dal progetto perch\'e i commenti degli scenari hanno affermato il
contrario per un certo periodo:

\begin{center}
\small
\begin{adjustbox}{max width=\textwidth}
\begin{tabular}{@{}lrr@{}}
\toprule
 & \textbf{\code{word-stats-java}} & \textbf{\code{word-stats-java-lite}} \\
\midrule
build & Spring Boot su JVM & binario nativo GraalVM \\
SDK & \code{sdks/java} & \code{sdks/java-lite} \\
residente, a riposo & 133\,MiB & 2{,}7\,MiB \\
residente, a caldo & 160\,MiB & 51{,}6\,MiB \\
tempo di servizio & 5{,}57\,ms & 2{,}15\,ms \\
\bottomrule
\end{tabular}
\end{adjustbox}
\end{center}

\begin{quote}
Qualunque esperimento che le tratti come intercambiabili \`e confuso.
\end{quote}

\`E riportato in questo capitolo, e non solo nel \cref{cap:18-sdk}, perch\'e
riguarda direttamente l'interpretazione della prima sezione: l'asimmetria fra le
due funzioni nella tabella delle concessioni \`e in larga parte questa
differenza di runtime, non una scelta dell'allocatore.

\section{Che cosa non \`e stabilito}

Il progetto enumera esplicitamente i limiti dei propri risultati. Sono riportati
qui integralmente perch\'e sono parte del contributo:

\begin{enumerate}[leftmargin=*]
  \item Perch\'e una raccolta completa costi 883\,ms quando l'insieme vivo \`e
  piccolo. La pausa traccia la dimensione dello heap anzich\'e i dati vivi, ma il
  meccanismo non \`e stato confermato.
  \item Se i pesi di \code{BUDGETED} facciano qualcosa di utile. Entrambe le
  funzioni hanno ricevuto SLO e pesi identici di proposito, quindi i livelli di
  servizio differenziati restano non verificati.
  \item Se il risultato di \code{SOJOURN} sotto arrivi aperti regga. \`E una sola
  esecuzione non ripetuta per modo, su un banco in cui entrambi i sistemi
  scartavano il 14--16\% del carico offerto.
  \item Che cosa faccia \code{SOJOURN} quando il sistema \emph{non} \`e
  sovraccarico. Ogni sua esecuzione finora \`e stata a capacit\`a o oltre, quindi
  il modo non \`e mai stato osservato scegliere un limite in condizioni comode.
  \item Se la soglia di guardia al 70\% sia corretta. \`E stata scelta, non
  misurata, e sul profilo a ciclo chiuso ha scavalcato completamente il modello.
\end{enumerate}

Il punto 2 merita enfasi: l'equit\`a max-min \emph{pesata} \`e presentata nel
\cref{cap:12-concurrency-control} come uno dei contributi del lavoro, e i pesi
non sono mai stati esercitati con valori diversi. Ci\`o che \`e stato validato \`e
l'allocazione a pesi uguali.

\section{Sintesi}

\begin{table}[htbp]
\centering
\caption{Le tre campagne e i loro esiti.}
\label{tab:campagne}
\small
\begin{tabularx}{\textwidth}{@{}lXX@{}}
\toprule
\textbf{Campagna} & \textbf{Risultato} & \textbf{Forza dell'evidenza} \\
\midrule
\code{BUDGETED} contro \code{ADAPTIVE\_PER\_POD} & Pari end-to-end; il primo con
met\`a della concorrenza sulla funzione che non ne ha bisogno, e totale sotto
budget. & Buona sugli aggregati; l'impianto a ciclo chiuso non poteva
distinguere i segnali. \\
\code{SOJOURN} contro \code{ADAPTIVE\_PER\_POD} & $+3{,}2\%$ servite,
$-15{,}7\%$ rifiuti, $-10{,}6\%$ p95, con \emph{meno} concorrenza. & Una sola
esecuzione per modo, entrambi in sovraccarico. \\
Co-tenancy & Cross-talk reale ($-28\%$ throughput), \emph{non} contesa di CPU
(differenza 0{,}6\%). & Solida: la manipolazione del \code{cpuset} \`e un
controllo diretto. \\
\bottomrule
\end{tabularx}
\end{table}

Il terzo risultato \`e il pi\`u solido dei tre, ed \`e negativo: esclude
un'ipotesi anzich\'e stabilirne una. \`E anche il pi\`u utile per il seguito del
lavoro, perch\'e restringe lo spazio di ricerca sul meccanismo del cross-talk al
control plane condiviso.
